\chapter{Самопорождающийся переходный дефект: минимальный нелинейный тест}

Язык первичных переходов становится физически содержательным только если
дефект не вставляется готовым профилем. Настоящая глава выполняет первый
минимальный тест: локальный функционал должен сам породить неоднородное
седло, защитить его топологическим числом, допустить перенос без изменения
формы и локализовать дираковскую нулевую моду.

Этот тест специально отделён от прежней $B-L$-ветви. Там профиль
$\tanh x$ не следовал из предъявленного действия. Здесь сначала задаётся
полный локальный функционал, затем профиль выводится из его уравнения и
воспроизводится независимой решёточной релаксацией.

\section{Литературная граница}

Самосогласованные неоднородные конденсаты и кинковые решения известны в
моделях Гросса--Невё; см. \path{arXiv:0806.2659}. Знакопеременный
дираковский массовый фон локализует нулевую моду по механизму
Джакива--Ребби; см. DOI \path{10.1103/PhysRevD.13.3398}. Нелинейные
квантовые блуждания могут сходиться к нелинейному уравнению Дирака;
см. \path{arXiv:1902.02017}. Взаимодействующие фермионные клеточные
автоматы также способны иметь связанные состояния; см.
\path{arXiv:2304.14687}.

Следовательно, новизной проекта не может быть само существование кинков,
нулевых мод или нелинейных автоматов. Проверяется только более узкое
соединение с носителем Мориты и проектными запретами.

\section{Минимальный локальный функционал}

В непрерывном пределе введём вещественное поле $q(x,t)$, управляющее
ориентационным смешиванием перехода. В безразмерных единицах рассмотрим
\begin{equation}
 S_q=\int dt\,dx\left[
 \frac12(\partial_tq)^2-\frac12(\partial_xq)^2
 -\frac14(q^2-1)^2
 \right].
 \label{eq:v5-self-defect-action}
\end{equation}
Уравнение поля равно
\begin{equation}
 \partial_t^2q-\partial_x^2q+q(q^2-1)=0.
 \label{eq:v5-self-defect-eom}
\end{equation}

В секторе граничных условий
\begin{equation}
 q(-\infty)=-1,
 \qquad
 q(+\infty)=+1
 \label{eq:v5-self-defect-boundary}
\end{equation}
статическое уравнение имеет решение
\begin{equation}
 q_*(x)=\tanh\frac{x-X}{\sqrt2}.
 \label{eq:v5-self-defect-kink}
\end{equation}
Здесь $X$ не является подогнанным профилем: это коллективная координата,
возникающая из трансляционной симметрии.

Подстановка в уравнение даёт нулевой остаток. Энергия решения равна
\begin{equation}
 E_{\rm kink}=\int dx\left[
 \frac12(q_*')^2+\frac14(q_*^2-1)^2
 \right]=\frac{2\sqrt2}{3}.
 \label{eq:v5-self-defect-energy}
\end{equation}
Топологический заряд
\begin{equation}
 Q=\frac{q(+\infty)-q(-\infty)}2=1
 \label{eq:v5-self-defect-charge}
\end{equation}
не меняется при локальной непрерывной деформации с сохранёнными границами.

\section{Проверка самопорождения профиля}

Для проверки, что формула \eqref{eq:v5-self-defect-kink} не была скрыто
вставлена, машинный аудит решает градиентный поток
\begin{equation}
 \partial_s q=\partial_x^2q-q(q^2-1)
 \label{eq:v5-self-defect-gradient-flow}
\end{equation}
на конечной решётке с границами $-1,+1$. Начальный профиль берётся
линейным с детерминированной поперечной возмущающей модой и не использует
$\tanh$. Энергия монотонно убывает, а результат сходится к аналитическому
кинку с малой среднеквадратичной ошибкой.

Это является положительным proof-of-concept: локальное нелинейное правило
способно породить профиль дефекта из уравнения, если топологический сектор
уже задан граничными условиями.

\section{Перенос без расплывания}

Действие \eqref{eq:v5-self-defect-action} лоренц-инвариантно. Поэтому для
$|v|<1$
\begin{equation}
 q_v(x,t)=
 \tanh\!\left(
 \frac{\gamma(x-vt-X)}{\sqrt2}
 \right),
 \qquad
 \gamma=(1-v^2)^{-1/2},
 \label{eq:v5-self-defect-moving}
\end{equation}
является точным движущимся решением. Оно сохраняет форму, заряд и
релятивистскую энергию $E(v)=\gamma E_{\rm kink}$.

Это впервые реализует слабую часть требуемого образа: перемещается не
одна и та же порция субстанции, а положение устойчивого дефектного класса.
Однако непрерывная лоренц-инвариантность здесь объявлена действием, а не
выведена из точного дискретного автомата.

\section{Дираковская нулевая мода}

Свяжем дефект с ориентационным дублетом минимального перехода:
\begin{equation}
 H_q=-i\sigma_2\partial_x+gq(x)\sigma_1.
 \label{eq:v5-self-defect-dirac}
\end{equation}
Для $g>0$ существует нормируемое нулевое решение
\begin{equation}
 \psi_0(x)=\mathcal N
 \begin{pmatrix}
 \cosh^{-g\sqrt2}(x/\sqrt2)\\0
 \end{pmatrix},
 \qquad H_{q_*}\psi_0=0.
 \label{eq:v5-self-defect-zero-mode}
\end{equation}
Его существование определяется сменой знака массы на бесконечностях, а не
точной формой кинка.

Но полный проектный перенос действует на
\begin{equation}
 \mathbb C^2\otimes E,
 \qquad E=M_{20\times15}(\mathbb C).
\end{equation}
Бимодульная ковариантность требует скалярного действия на $E$. Поэтому
оператор $H_q\otimes I_E$ локализует не одну, а
\begin{equation}
 \dim_{\mathbb C}E=300
 \end{equation}
одинаковых нулевых мод. Минимальная нелинейность создаёт дефект, но не
выбирает наблюдаемый сектор.

\begin{nogocriterion}[Запрет единственной частицы на полном носителе]
Скалярный знакопеременный массовый дефект на полном соответствии
$M_{20\times15}$ наследует всю внутреннюю кратность $300$. Получение одной
или малого физического числа нулевых мод требует дополнительного
проектора, нескалярной связности либо индексного комплекса высшей степени.
Следовательно, один универсальный кинк не является спектром частиц.
\end{nogocriterion}

\section{Цена положительного примера}

В размерной форме функционал содержит как минимум
\begin{equation}
 V(q)=\frac\lambda4(q^2-v_q^2)^2,
 \qquad
 H_q=-i\sigma_2\partial_x+yq\sigma_1.
 \label{eq:v5-self-defect-dimensional-data}
\end{equation}
После выбора единиц остаются физический масштаб и безразмерное отношение
типа $y/\sqrt\lambda$. Ни двойная яма, ни её масштаб, ни это отношение не
выведены из $M_{35}$, следа $4/7,3/7$ или четырёхтактной фазы.

Кроме того, заряд $Q=1$ не возникает из тривиальных одинаковых границ:
топологический сектор задаётся условиями
\eqref{eq:v5-self-defect-boundary}. Динамика выводит профиль внутри сектора,
но не объясняет, почему Вселенная выбрала ненулевой суммарный сектор.

\begin{theorem}[Положительный механизм и отрицательное замыкание]
Минимальный локальный нелинейный функционал действительно сам выводит
устойчивый профиль, его движущееся продолжение и локализованную дираковскую
нулевую моду. Тем самым он доказывает динамическую непротиворечивость идеи
частицы-дефекта. Но на полном носителе Мориты возникает кратность $300$, а
потенциал, масштаб, юкавское отношение и топологические границы являются
новыми входами. Физическая теория частиц не замкнута.
\end{theorem}

\section{Статус следующего шага}

Продолжать подбором более сложного потенциала запрещено. Единственный
содержательный вопрос после этого proof-of-concept состоит в том, может ли
уже существующий проектор голономной нулевой линии войти в нелинейное
правило функториально и сократить кратность $300$ без ручного выбора
частицы. Если нет, часть II должна завершиться: общий механизм дефекта
возможен, но его физическое внутреннее представление текущим родителем не
выведено.

\begin{protocol}
\begin{enumerate}
 \item локальный нелинейный функционал: \textbf{предъявлен};
 \item профиль $\tanh$: \textbf{выведен, а не вставлен};
 \item независимая решёточная релаксация: \textbf{выполнена};
 \item топологический заряд: \textbf{$Q=1$};
 \item перенос без расплывания: \textbf{получен в непрерывной модели};
 \item нормируемая дираковская нулевая мода: \textbf{получена};
 \item точный дискретный нелинейный автомат: \textbf{не построен};
 \item кратность на полном носителе: \textbf{$300$};
 \item потенциал и относительная связь из родителя: \textbf{не выведены};
 \item физическое замыкание: \textbf{не достигнуто}.
\end{enumerate}
\end{protocol}

Машинный протокол сохранён в
\path{s2t_v5_self_generated_transition_defect_gate_results.json}.